\section{Overview of Inverse Compton Scattering EUV Light Source}

\subsection{Physics of Inverse Compton Scattering}

Inverse Compton scattering (ICS) describes photon scattering from relativistic
electrons. In the laboratory frame, the scattered energy and angular distribution
depend on electron energy, incident-photon energy, collision and observation
geometry, polarization, and recoil assumptions~\cite{SunWu2011ComptonSource}.
For intense pulsed fields, nonlinear motion and pulse geometry can modify the
spectrum; relevant treatments include general nonlinear Thomson geometries and
ponderomotive broadening~\cite{KrafftEtAl2005NonlinearThomson,HartemannWu2013Brightness}.
The EUVICS calculation must state which regime and conventions it uses rather
than combining results from these models without matched assumptions.

\subsection{Accelerator System}

The candidate system boundary includes an electron source, injector, energy-setting
system, transport, focusing, steering, diagnostics, and dump. The specific
technology, beam parameters, interfaces, and verification criteria remain TBD.
Published SRF photoinjector reviews and experiments provide evidence about other
machines, not EUVICS performance~\cite{ArnoldTeichert2011SRFPhotoinjectors,PetrushinaEtAl2020SRFGun}.

\subsection{Laser System}

The candidate laser boundary includes pulse generation, transport, focusing,
polarization control, synchronization, diagnostics, and the interaction-region
interface. Wavelength, pulse energy, duration, waist, bandwidth, polarization,
timing, and stability values remain TBD except for the roadmap's 800~nm nominal
assumed input. Laser-to-RF synchronization and time-stamping methods are reviewed
in the ultrafast-electron literature~\cite{FilippettoEtAl2022UED}; EUVICS-specific
tolerances require an approved overlap and performance budget.

\blankline
\blankline
